\subsection{问题三模型的建立与求解}


\subsubsection{多故障源信号数据预处理}


在多故障源情形下，观测信号中可能同时包含多个不同频率的微弱周期分量。设附件中“多故障源”sheet 给出的原始观测序列为


\[
x_i=x(t_i),\qquad i=1,2,\cdots,N.
\]



与前两问保持一致，首先对原始信号进行去均值处理：


\[
y_i=x_i-\bar{x},\qquad 
\bar{x}=\frac{1}{N}\sum_{i=1}^{N}x_i.
\]


由于本问仍需估计各故障分量的振幅，主模型保留去均值信号的原始幅值尺度，不再作幅值归一化。


多故障源信号模型写为


\[
y(t_i)=\sum_{k=1}^{K}A_k\sin(2\pi f_k t_i+\varphi_k)+\varepsilon_i,
\]


其中 \(K\) 为故障源个数，\(f_k,A_k,\varphi_k\) 分别为第 \(k\) 个故障源的特征频率、振幅和初相，\(\varepsilon_i\) 表示噪声及未建模扰动。




\subsubsection{基于全长正弦投影-OMP与局部变量投影精修的多故障源分离模型}


本节将前两问的单频投影和最小二乘恢复推广到多频情形，建立两阶段分离模型：先用全长正弦投影-OMP 在残差中自动寻找多个显著频率，再用局部变量投影精修近频或宽峰区域。


\paragraph{（1）多正弦模型的线性化表示}


沿用前两问的线性化思想，对任意候选频率 \(f\)，单个正弦分量写为


\[
A\sin(2\pi ft+\varphi)
=a\cos(2\pi ft)+b\sin(2\pi ft),
\]


其中


\[
a=A\sin\varphi,\qquad b=A\cos\varphi.
\]


于是对候选频率 \(f\)，构造正弦投影矩阵


\[
\mathbf H(f)=
\begin{bmatrix}
\cos(2\pi f t_1)&\sin(2\pi f t_1)\\
\cos(2\pi f t_2)&\sin(2\pi f t_2)\\
\vdots&\vdots\\
\cos(2\pi f t_N)&\sin(2\pi f t_N)
\end{bmatrix}.
\]


若频率 \(f\) 给定，余弦、正弦系数可由最小二乘直接估计：


\[
\hat{\boldsymbol\theta}(f)
=
\left[\mathbf H^T(f)\mathbf H(f)\right]^{-1}
\mathbf H^T(f)\mathbf y.
\]


对应的拟合分量为


\[
\hat{\mathbf s}(f)=\mathbf H(f)\hat{\boldsymbol\theta}(f).
\]


当 \(f\) 接近真实故障频率时，观测信号在该二维正弦子空间上的投影能量明显增大，因此可沿用问题一的投影能量作为多源搜索的定位函数。


\paragraph{（2）基于残差投影的 OMP 频率搜索}


第 \(m\) 次迭代前，设当前残差为 \(\mathbf r^{(m-1)}\)。对任意候选频率 \(f\)，定义残差投影能量


\[
J_m(f)=
{\mathbf r^{(m-1)}}^{T}\mathbf H(f)
\left[\mathbf H^{T}(f)\mathbf H(f)\right]^{-1}
\mathbf H^{T}(f)\mathbf r^{(m-1)}.
\]


第 \(m\) 个故障候选频率由


\[
\hat f_m=\arg\max_f J_m(f)
\]


确定。该过程等价于在连续正弦字典中寻找与当前残差最相关的频率分量，使强分量被解释后，较弱故障峰能在后续残差投影谱中显现。为避免逐个扣除造成误差累积，每识别出一个新频率后，本文对当前所有已识别频率进行多频联合最小二乘重拟合。设已识别频率集合为


\[
\mathcal F_m=\{\hat f_1,\hat f_2,\cdots,\hat f_m\}.
\]


构造联合矩阵


\[
\mathbf H_m=
\left[
\cos(2\pi \hat f_1\mathbf t),\sin(2\pi \hat f_1\mathbf t),
\cdots,
\cos(2\pi \hat f_m\mathbf t),\sin(2\pi \hat f_m\mathbf t)
\right].
\]


联合系数估计为


\[
\hat{\boldsymbol\theta}_m
=
(\mathbf H_m^T\mathbf H_m)^{-1}\mathbf H_m^T\mathbf y,
\]


重构信号和残差分别为


\[
\hat{\mathbf s}_m=\mathbf H_m\hat{\boldsymbol\theta}_m,
\qquad
\mathbf r^{(m)}=\mathbf y-\hat{\mathbf s}_m.
\]


联合最小二乘使多个分量共同解释原始信号，减少识别顺序对振幅、相位和残差的影响。


最后，对第 \(k\) 个频率，其余弦、正弦系数记为 \(a_k,b_k\)，则振幅和初相恢复为


\[
\hat A_k=\sqrt{a_k^2+b_k^2},
\qquad
\hat\varphi_k=\arctan2(a_k,b_k).
\]


\paragraph{（3）局部变量投影精修}


当两个故障频率非常接近时，单频投影峰可能重叠，OMP 可能得到合并峰。为提高近频分辨能力，本文引入局部双频变量投影精修。


变量投影法将多频估计写成可分离非线性最小二乘问题。对给定频率向量


\[
\mathbf f=[f_1,f_2,\cdots,f_K]^T,
\]


构造多频字典矩阵


\[
\boldsymbol\Psi(\mathbf f)=
\left[
\cos(2\pi f_1\mathbf t),\sin(2\pi f_1\mathbf t),
\cdots,
\cos(2\pi f_K\mathbf t),\sin(2\pi f_K\mathbf t)
\right].
\]


多频拟合问题为


\[
\min_{\mathbf f,\mathbf c}
\left\|
\mathbf y-\boldsymbol\Psi(\mathbf f)\mathbf c
\right\|_2^2.
\]


其中 \(\mathbf f\) 是非线性频率参数，\(\mathbf c\) 是线性的振幅相位系数。对任意固定的 \(\mathbf f\)，线性系数都有闭式最小二乘解


\[
\hat{\mathbf c}(\mathbf f)
=
\boldsymbol\Psi^\dagger(\mathbf f)\mathbf y,
\]


其中 \(\boldsymbol\Psi^\dagger(\mathbf f)\) 表示 Moore--Penrose 伪逆。令


\[
P_{\boldsymbol\Psi(\mathbf f)}
=
\boldsymbol\Psi(\mathbf f)\boldsymbol\Psi^\dagger(\mathbf f)
\]


为投影到多频子空间的正交投影矩阵，则原问题可约化为只关于频率的优化问题：


\[
\min_{\mathbf f}
\left\|
\left(I-P_{\boldsymbol\Psi(\mathbf f)}\right)\mathbf y
\right\|_2^2.
\]


这样可先消去线性的振幅相位系数，只对频率向量 \(\mathbf f\) 进行优化，降低近频精修的搜索维数。


在本文实现中，变量投影作为 OMP 后的局部增强步骤。当某一局部频带可能包含两个相近频率时，以 OMP 或全长投影给出的峰值位置作为初值中心，构造双频模型


\[
\min_{f_1,f_2}
\left\|
\mathbf y-\boldsymbol\Psi(f_1,f_2)
\hat{\mathbf c}(f_1,f_2)
\right\|_2^2,
\]


其中


\[
\hat{\mathbf c}(f_1,f_2)
=
\left[\boldsymbol\Psi^T(f_1,f_2)\boldsymbol\Psi(f_1,f_2)\right]^{-1}
\boldsymbol\Psi^T(f_1,f_2)\mathbf y.
\]


然后在局部频率区间内搜索使约化残差最小的 \((f\_1,f\_2)\)。该步骤保留 OMP 的自动性，并用于突破普通 FFT 的频率栅格限制。


\subsubsection{故障源数判定}


为了避免故障源个数误判，本文根据残差投影谱的稳健噪声底自动判定是否继续加入新的故障频率。若某一峰值仅处于背景波动范围内，则不将其解释为新的故障源。


首先，在每次残差投影谱中排除已识别频率附近的邻域，避免同一频率峰被重复选择。本文设置最小频率间隔为


\[
\Delta f_{\min}=0.01\ \mathrm{Hz}.
\]


然后，对有效频率区间上的残差投影能量样本 \(\{J\_m(f\_j)\}\) 计算中位数


\[
M_m=\mathrm{median}\{J_m(f_j)\},
\]


以及 MAD 稳健尺度估计


\[
\sigma_{\mathrm{MAD},m}
=1.4826\cdot
\mathrm{median}\left(|J_m(f_j)-M_m|\right).
\]


其中 \(1.4826\) 为 MAD 一致性修正系数。相比均值和标准差，中位数与 MAD 对少数强故障峰不敏感，更适合强峰、弱峰与噪声共存时的源数判定。


停止阈值定义为


\[
\eta_m=M_m+\lambda\sigma_{\mathrm{MAD},m},
\]


其中本文取


\[
\lambda=18.
\]


该取值较保守，可降低随机噪声峰被误判为故障频率的风险。


即当前最大峰值 \(J_m^{\max}\) 同时满足


\[
J_m^{\max}\ge \eta_m
\]


以及


\[
\frac{J_m^{\max}}{M_m}\ge 20
\]


时，才认为残差中仍存在显著周期故障分量，并继续加入该频率。否则停止迭代，并将当前已识别频率个数作为故障源数估计


\[
\hat K=m-1.
\]


峰值/噪声底比值阈值 20 是对 \(\eta_m\) 的补充约束，用于避免残差能量较低时由小随机起伏触发误检。结合最小频率间隔 \(\Delta f_{\min}=0.01\ \mathrm{Hz}\)，可防止同一宽峰或旁瓣被重复计为多个故障源。


\subsubsection{问题求解过程}


对 data.xlsx 中“多故障源”数据进行计算。图中给出了原始多故障信号对应的全长正弦投影能量谱 \(J\_1(f)\)，以及自动识别出的多个故障特征频率。


\begin{figure}[H]
\centering
\includegraphics[width=0.86\textwidth]{figures/多故障源分离_run_q3_multisource_separation/fig_q3_original_spectrum_identified.png}
\caption{全长正弦投影能量谱与识别频率}
\end{figure}


投影-OMP 的迭代过程如下表所示。


\begin{table}[H]
\centering
\caption{OMP 迭代识别过程}
\small
\begin{tabular}{cccccc}
\toprule
迭代 & 粗定位频率 / Hz & 精估频率 / Hz & 峰值投影能量 & 峰值/噪声底 & 累计能量占比 \\
\midrule
1 & 13.999650 & 14.000024 & 29.796625 & 2163.361 & 0.070822 \\
2 & 7.999800 & 7.999929 & 12.279573 & 892.686 & 0.098171 \\
3 & 3.999900 & 3.999998 & 7.263044 & 528.815 & 0.114281 \\
4 & 12.999675 & 13.000089 & 2.028304 & 147.679 & 0.119152 \\
\bottomrule
\end{tabular}
\end{table}


第 5 次迭代时，残差投影谱最大峰值约为


\[
0.185084,
\]


而稳健阈值约为


\[
0.270040,
\]


且峰值/噪声底比值约为


\[
13.476,
\]


低于设定阈值 20。停止继续扫描故障源，得到故障源个数


\[
\hat K=4.
\]


残差投影谱的逐步变化如下图所示。可以看到，强分量被依次提取后，对应频率峰在后续残差中明显削弱。


\begin{figure}[H]
\centering
\includegraphics[width=0.86\textwidth]{figures/多故障源分离_run_q3_multisource_separation/fig_q3_omp_residual_projection_spectra.png}
\caption{OMP逐步残差投影谱}
\end{figure}


最终识别出的故障频率、振幅和初相如下表所示。


\begin{table}[H]
\centering
\caption{多故障源参数估计结果}
\small
\begin{tabular}{ccccc}
\toprule
故障源 & 频率 / Hz & 振幅 & 初相 / rad & 能量占比 \\
\midrule
1 & 3.999998185 & 0.019102914 & 1.371543820 & 0.016111 \\
2 & 7.999928733 & 0.024889374 & -2.989236395 & 0.027349 \\
3 & 13.000088597 & 0.010502940 & -2.157357956 & 0.004870 \\
4 & 14.000024160 & 0.040052738 & 0.111605514 & 0.070822 \\
\bottomrule
\end{tabular}
\end{table}


因此，多故障源信号可恢复为


\[
\hat s(t)=\sum_{k=1}^{4}
\hat A_k\sin(2\pi \hat f_k t+\hat\varphi_k).
\]


代入数值可写为


\[
\begin{aligned}
\hat s(t)=&
0.019102914\sin(2\pi\cdot3.999998185t+1.371543820)\\
&+0.024889374\sin(2\pi\cdot7.999928733t-2.989236395)\\
&+0.010502940\sin(2\pi\cdot13.000088597t-2.157357956)\\
&+0.040052738\sin(2\pi\cdot14.000024160t+0.111605514).
\end{aligned}
\]


多频重构后的解释能量占比为


\[
R^2=0.119151588464,
\]


残差 RMS 为


\[
0.099879873631,
\]


相对于原始去均值信号的 RMS 降低约为


\[
0.550988\ \mathrm{dB}.
\]


下图展示了短时段内观测信号、多频重构信号和残差。故障分量相对于噪声较弱，因此重构信号不会完全贴合原始波形，但其周期结构已通过全长相干积累被提取出来。


\begin{figure}[H]
\centering
\includegraphics[width=0.86\textwidth]{figures/多故障源分离_run_q3_multisource_separation/fig_q3_time_reconstruction_short.png}
\caption{多频重构与残差}
\end{figure}




频域对比表明，重构信号集中在识别出的 4 个故障频率附近，而残差中这些主峰明显削弱。


\begin{figure}[H]
\centering
\includegraphics[width=0.86\textwidth]{figures/多故障源分离_run_q3_multisource_separation/fig_q3_frequency_before_after.png}
\caption{多源分离前后频域对比}
\end{figure}


为检验源数判定的稳定性，改变稳健阈值倍数 \(\lambda\)：


\[
\lambda=12,14,16,18,20,24.
\]


计算结果均自动识别出 4 个故障源，频率列表稳定为约 \(4,8,13,14\ \mathrm{Hz}\)，重构能量占比均为 \(0.119151588464\)，残差 RMS 均为 \(0.099879873631\)，说明源数判定对阈值倍数不敏感。


\subsubsection{近频辨识极限实验}


为考察频率接近时的辨识极限，构造双频仿真实验：


\[
x(t)=A_1\sin(2\pi f_1t+\varphi_1)
+A_2\sin(2\pi f_2t+\varphi_2)+n(t),
\]


其中 \(f_1=10\ \mathrm{Hz}\)，\(f_2=f_1+\Delta f\)。因此，在每一组仿真中真实频率为


\[
\mathbf f_{\mathrm{true}}=
\left[
10,\quad 10+\Delta f
\right]^T\ \mathrm{Hz}.
\]


采样频率和观测时长与附件数据保持一致，即 \(f_s=100\ \mathrm{Hz}\)、\(T=400\ \mathrm{s}\)。改变频率间隔 \(\Delta f\) 和信噪比 SNR，多次重复实验并统计两个频率均被正确识别的比例。


本文将“正确识别”定义为：算法输出的估计频率中，至少存在两个不同频率 \(\hat f_a,\hat f_b\)，分别对应两个真实频率，且误差均不超过容许范围 \(\varepsilon\)，即


\[
\begin{aligned}
|\hat f_a-f_1|&\le \varepsilon,\qquad \\
|\hat f_b-f_2|&\le \varepsilon,\qquad a\ne b.
\end{aligned}
\]


其中变量投影对比实验采用的误差容限为


\[
\varepsilon=\max\left(0.00012,\ \min(0.0012,\ 0.35\Delta f)\right)\ \mathrm{Hz}.
\]


该判据要求两个真实频率分别由两个不同估计频率匹配；若只给出合并峰或误差超限，则记为识别失败。




\begin{figure}[H]
\centering
\includegraphics[width=0.86\textwidth]{figures/近频超分辨实验_run_q3_superresolution_experiment/fig_q3_superresolution_fft_vs_varpro.png}
\caption{FFT与变量投影近频分辨率对比}
\end{figure}


对比结果显示，传统 FFT 双峰拾取在 \(\Delta f<0.005\ \mathrm{Hz}\) 时基本不能稳定分辨两个近频分量；变量投影法在较高信噪比下可实现低于传统 FFT 栅格间隔的分辨能力。例如，\(\mathrm{SNR}=-8\ \mathrm{dB}\) 且 \(\Delta f=0.001\ \mathrm{Hz}\) 时识别成功率达到 1；\(\mathrm{SNR}=-20\ \mathrm{dB}\) 且 \(\Delta f=0.002\ \mathrm{Hz}\) 时识别成功率也达到 1。说明全长投影-OMP 的频率初值可通过局部双频变量投影进一步精修。


\subsubsection{本节小结}


本节将前两问的全长正弦投影模型推广到多故障源情形，建立了“全长正弦投影-OMP + 局部变量投影精修”的分离模型。该模型识别出 4 个故障源，频率约为 \(4,8,13,14\ \mathrm{Hz}\)，并给出对应振幅和初相。近频实验表明，变量投影精修可提升低于传统 FFT 栅格间隔的近频分辨能力。
